% GENERATED FILE. Edit canonical lessons, then run npm run generate:books.
% canonical-source-hash: fnv1a64:5363ef2dafcb4279
% canonical-lessons: PA-C343-ulajhna, PA-C343-ukkna, PA-C343-ishara, PA-C343-mamla, PA-C343-sabut, PA-R343-review-337-343

\chapter{A Mix-up}
\label{ch:pa-a2-a-mix-up}
\hlchaptermodality{343}

\begin{chapteropening}
\textbf{By the end of this chapter:} \emph{I can say that I got muddled and sort out a mix-up.}

Complete the last lesson of chapter 343: I can say that I got muddled and sort out a mix-up.
\end{chapteropening}

\section[\texorpdfstring{\pa{ਉਲਝਣਾ} (ulajhṇā)}{ਉਲਝਣਾ (ulajhṇā)}]{\pa{ਉਲਝਣਾ} (ulajhṇā) — to get tangled; to get muddled}

\label{lesson:PA-C343-ulajhna}



\begin{quote}
Before the new one: say the Punjabi for plain to understand, then the Punjabi for to suggest.
\end{quote}

\subsection*{You'll want to know: \pa{ਉਲਝਣਾ}}
\textbf{\pa{ਉਲਝਣਾ}} — \emph{ulajhṇā} — \textquotedblleft{}to get tangled; to get muddled\textquotedblright{}.

\begin{quote}
\pa{ਮੈਂ} \pa{ਲੰਮੇ} \pa{ਸਵਾਲਾਂ} \pa{ਵਿੱਚ} \pa{ਨਹੀਂ} \pa{ਉਲਝਦਾ} \pa{ਹਾਂ।}
\end{quote}

\begin{quote}
\emph{maiṁ lamme savālāṁ vich nahīṁ ulajhdā hāṁ.}
\end{quote}

\begin{quote}
I don't get muddled by long questions.
\end{quote}

Past (I …): \pa{ਮੈਂ} \pa{ਸਵਾਲ} \pa{ਵਿੱਚ} \pa{ਉਲਝਿਆ}, \emph{maiṁ savāl vich uljhiā}.

\begin{cousinweb}
\textbf{\pa{ਉਲਝਣਾ}} (\emph{ulajhṇā}) is inherited from Sanskrit \emph{uparundhati}, \textquotedblleft{}obstructs\textquotedblright{}; CDIAL lists Punjabi \emph{uljhaṇā}, \textquotedblleft{}to be entangled\textquotedblright{}, under it (CDIAL 2221).
\end{cousinweb}

\begin{grammarlens}[title={What you've built}]
One more word for asking and answering.
\end{grammarlens}

\subsection*{Your turn}
\begin{itemize}
\raggedright
  \item \emph{Say it:} \emph{ulajhṇā}
  \item \emph{Say it:} \emph{ulajhṇā}, once more
  \item \emph{Say it:} \emph{maiṁ lamme savālāṁ vich nahīṁ ulajhdā hāṁ.} — I don't get muddled by long questions.
  \item \emph{Recall:} say the Punjabi for plain to understand, then the Punjabi for to suggest, then say \emph{ulajhṇā} again
\end{itemize}

\begin{tcolorbox}[breakable,colback=teal!4,colframe=teal!35!black,title={Before you move on}]
What does \pa{ਉਲਝਣਾ} mean? (\textquotedblleft{}To get tangled; to get muddled\textquotedblright{}.) Say it once more.
\end{tcolorbox}
\section[\texorpdfstring{\pa{ਉੱਕਣਾ} (ukkṇā)}{ਉੱਕਣਾ (ukkṇā)}]{\pa{ਉੱਕਣਾ} (ukkṇā) — to slip up; to make a mistake}

\label{lesson:PA-C343-ukkna}



\begin{quote}
Before the new one: say the Punjabi for to suggest, then the Punjabi for to get muddled.
\end{quote}

\subsection*{You'll want to know: \pa{ਉੱਕਣਾ}}
\textbf{\pa{ਉੱਕਣਾ}} — \emph{ukkṇā} — \textquotedblleft{}to slip up; to make a mistake\textquotedblright{}.

\begin{quote}
\pa{ਅਮਨ} \pa{ਹਿਸਾਬ} \pa{ਵਿੱਚ} \pa{ਕਦੇ} \pa{ਨਹੀਂ} \pa{ਉੱਕਦਾ} \pa{ਹੈ।}
\end{quote}

\begin{quote}
\emph{Aman hisāb vich kade nahīṁ ukkdā hai.}
\end{quote}

\begin{quote}
Aman never slips up in his sums.
\end{quote}

Past (I …): \pa{ਮੈਂ} \pa{ਜਵਾਬ} \pa{ਵਿੱਚ} \pa{ਉੱਕਿਆ}, \emph{maiṁ javāb vich ukkiā}.

\begin{cousinweb}
\textbf{\pa{ਉੱਕਣਾ}} (\emph{ukkṇā}) is inherited from Sanskrit \emph{utkramati}, \textquotedblleft{}ascends\textquotedblright{}, later \textquotedblleft{}neglects, transgresses\textquotedblright{}; CDIAL lists Punjabi \emph{ukkṇā}, \textquotedblleft{}to make a mistake\textquotedblright{}, under it (CDIAL 1737).
\end{cousinweb}

\begin{grammarlens}[title={What you've built}]
One more word for asking and answering.
\end{grammarlens}

\subsection*{Your turn}
\begin{itemize}
\raggedright
  \item \emph{Say it:} \emph{ukkṇā}
  \item \emph{Say it:} \emph{ukkṇā}, once more
  \item \emph{Say it:} \emph{Aman hisāb vich kade nahīṁ ukkdā hai.} — Aman never slips up in his sums.
  \item \emph{Recall:} say the Punjabi for to suggest, then the Punjabi for to get muddled, then say \emph{ukkṇā} again
  \item \emph{Write it:} copy \textbf{\pa{ਉੱਕਣਾ}} once, checking the whole word
\end{itemize}

\begin{tcolorbox}[breakable,colback=teal!4,colframe=teal!35!black,title={Before you move on}]
What does \pa{ਉੱਕਣਾ} mean? (\textquotedblleft{}To slip up; to make a mistake\textquotedblright{}.) Say it once more.
\end{tcolorbox}
\section[\texorpdfstring{\pa{ਇਸ਼ਾਰਾ} (ishārā)}{ਇਸ਼ਾਰਾ (ishārā)}]{\pa{ਇਸ਼ਾਰਾ} (ishārā) — a hint; a sign made to someone}

\label{lesson:PA-C343-ishara}



\begin{quote}
Before the new one: say the Punjabi for to get muddled, then the Punjabi for to slip up.
\end{quote}

\subsection*{You'll want to know: \pa{ਇਸ਼ਾਰਾ}}
\textbf{\pa{ਇਸ਼ਾਰਾ}} — \emph{ishārā} — \textquotedblleft{}a hint; a sign made to someone\textquotedblright{}.

\begin{quote}
\pa{ਕੀ} \pa{ਤੁਸੀਂ} \pa{ਮੇਰਾ} \pa{ਇਸ਼ਾਰਾ} \pa{ਸਮਝਦੇ} \pa{ਹੋ}?
\end{quote}

\begin{quote}
\emph{kī tusī̃ merā ishārā samajhde ho?}
\end{quote}

\begin{quote}
Do you get my hint?
\end{quote}

A masculine noun.

\begin{cousinweb}
\textbf{\pa{ਇਸ਼ਾਰਾ}} (\emph{ishārā}) is a loan, usually traced through Persian to Arabic \emph{ishāra}, \textquotedblleft{}sign\textquotedblright{}.
\end{cousinweb}

\begin{grammarlens}[title={What you've built}]
One more word for asking and answering.
\end{grammarlens}

\subsection*{Your turn}
\begin{itemize}
\raggedright
  \item \emph{Say it:} \emph{ishārā}
  \item \emph{Say it:} \emph{ishārā}, once more
  \item \emph{Say it:} \emph{kī tusī̃ merā ishārā samajhde ho?} — Do you get my hint?
  \item \emph{Recall:} say the Punjabi for to get muddled, then the Punjabi for to slip up, then say \emph{ishārā} again
  \item \emph{Write it:} copy \textbf{\pa{ਇਸ਼ਾਰਾ}} once, checking the whole word
\end{itemize}

\begin{tcolorbox}[breakable,colback=teal!4,colframe=teal!35!black,title={Before you move on}]
What does \pa{ਇਸ਼ਾਰਾ} mean? (\textquotedblleft{}A hint; a sign made to someone\textquotedblright{}.) Say it once more.
\end{tcolorbox}
\section[\texorpdfstring{\pa{ਮਾਮਲਾ} (māmlā)}{ਮਾਮਲਾ (māmlā)}]{\pa{ਮਾਮਲਾ} (māmlā) — a matter; an affair}

\label{lesson:PA-C343-mamla}



\begin{quote}
Before the new one: say the Punjabi for to slip up, then the Punjabi for a hint.
\end{quote}

\subsection*{You'll want to know: \pa{ਮਾਮਲਾ}}
\textbf{\pa{ਮਾਮਲਾ}} — \emph{māmlā} — \textquotedblleft{}a matter; an affair\textquotedblright{}.

\begin{quote}
\pa{ਇਹ} \pa{ਮਾਮਲਾ} \pa{ਸੌਖਾ} \pa{ਨਹੀਂ} \pa{ਹੈ।}
\end{quote}

\begin{quote}
\emph{eh māmlā saukhā nahīṁ hai.}
\end{quote}

\begin{quote}
This matter is not simple.
\end{quote}

A masculine noun.

\begin{cousinweb}
\textbf{\pa{ਮਾਮਲਾ}} (\emph{māmlā}) is a loan, usually traced through Persian to Arabic \emph{muʿāmala}, \textquotedblleft{}dealing\textquotedblright{}.
\end{cousinweb}

\begin{grammarlens}[title={What you've built}]
One more word for asking and answering.
\end{grammarlens}

\subsection*{Your turn}
\begin{itemize}
\raggedright
  \item \emph{Say it:} \emph{māmlā}
  \item \emph{Say it:} \emph{māmlā}, once more
  \item \emph{Say it:} \emph{eh māmlā saukhā nahīṁ hai.} — This matter is not simple.
  \item \emph{Recall:} say the Punjabi for to slip up, then the Punjabi for a hint, then say \emph{māmlā} again
  \item \emph{Write it:} copy \textbf{\pa{ਮਾਮਲਾ}} once, checking the whole word
\end{itemize}

\begin{tcolorbox}[breakable,colback=teal!4,colframe=teal!35!black,title={Before you move on}]
What does \pa{ਮਾਮਲਾ} mean? (\textquotedblleft{}A matter; an affair\textquotedblright{}.) Say it once more.
\end{tcolorbox}
\section[\texorpdfstring{\pa{ਸਬੂਤ} (sabūt)}{ਸਬੂਤ (sabūt)}]{\pa{ਸਬੂਤ} (sabūt) — proof; evidence}

\label{lesson:PA-C343-sabut}



\begin{quote}
Before the new one: say the Punjabi for a hint, then the Punjabi for a matter.
\end{quote}

\subsection*{You'll want to know: \pa{ਸਬੂਤ}}
\textbf{\pa{ਸਬੂਤ}} — \emph{sabūt} — \textquotedblleft{}proof; evidence\textquotedblright{}.

\begin{quote}
\pa{ਕੀ} \pa{ਤੁਹਾਡੇ} \pa{ਕੋਲ} \pa{ਕੋਈ} \pa{ਸਬੂਤ} \pa{ਹੈ}?
\end{quote}

\begin{quote}
\emph{kī tuhāḍe kol koī sabūt hai?}
\end{quote}

\begin{quote}
Do you have any proof?
\end{quote}

A masculine noun.

\begin{cousinweb}
\textbf{\pa{ਸਬੂਤ}} (\emph{sabūt}) is a loan, usually traced through Persian to Arabic \emph{thubūt}, \textquotedblleft{}firmness\textquotedblright{}.
\end{cousinweb}

\begin{grammarlens}[title={What you've built}]
One more word for asking and answering.
\end{grammarlens}

\subsection*{Your turn}
\begin{itemize}
\raggedright
  \item \emph{Say it:} \emph{sabūt}
  \item \emph{Say it:} \emph{sabūt}, once more
  \item \emph{Say it:} \emph{kī tuhāḍe kol koī sabūt hai?} — Do you have any proof?
  \item \emph{Recall:} say the Punjabi for a hint, then the Punjabi for a matter, then say \emph{sabūt} again
  \item \emph{Write it:} copy \textbf{\pa{ਸਬੂਤ}} once, checking the whole word
\end{itemize}

\begin{tcolorbox}[breakable,colback=teal!4,colframe=teal!35!black,title={Before you move on}]
What does \pa{ਸਬੂਤ} mean? (\textquotedblleft{}Proof; evidence\textquotedblright{}.) Say it once more.
\end{tcolorbox}
\section[(dialogue)]{Looking back: from Asking Nicely to A Mix-up}

\label{lesson:PA-R343-review-337-343}



\begin{quote}
Say the Punjabi for to ask for and a request.
\end{quote}

\subsection*{Your turn}
Say these aloud:

\begin{itemize}
\raggedright
  \item to ask for, a request, a little, a bit (to soften a request), to send for
  \item permission, forbidden, approved, an order, to interrupt
  \item a refusal, to back out, an excuse, by no means, to dig one's heels in
  \item a promise, to persuade, absolutely, to admit, a compromise
  \item doubt, certainty, trust, a guess, to seem right
  \item to explain, meaning, a reason, plain to understand, to suggest
  \item to get muddled, to slip up, a hint, a matter, proof
\end{itemize}

\begin{tcolorbox}[breakable,colback=teal!4,colframe=teal!35!black,title={Before you move on}]
To ask for? (\textbf{\pa{ਮੰਗਣਾ}}, \emph{maṅgṇā}.) A request? (\textbf{\pa{ਬੇਨਤੀ}}, \emph{bentī}.) A little? (\textbf{\pa{ਥੋੜ੍ਹਾ}}, \emph{thoṛhā}.) A bit (to soften a request)? (\textbf{\pa{ਜ਼ਰਾ}}, \emph{zarā}.) To send for? (\textbf{\pa{ਮੰਗਵਾਉਣਾ}}, \emph{maṅgvāuṇā}.) Permission? (\textbf{\pa{ਇਜਾਜ਼ਤ}}, \emph{ijāzat}.) Forbidden? (\textbf{\pa{ਮਨ੍ਹਾ}}, \emph{manhā}.) Approved? (\textbf{\pa{ਮਨਜ਼ੂਰ}}, \emph{manzūr}.) An order? (\textbf{\pa{ਹੁਕਮ}}, \emph{hukam}.) To interrupt? (\textbf{\pa{ਟੋਕਣਾ}}, \emph{ṭokṇā}.) A refusal? (\textbf{\pa{ਇਨਕਾਰ}}, \emph{inkār}.) To back out? (\textbf{\pa{ਟਲਣਾ}}, \emph{ṭalṇā}.) An excuse? (\textbf{\pa{ਬਹਾਨਾ}}, \emph{bahānā}.) By no means? (\textbf{\pa{ਹਰਗਿਜ਼}}, \emph{hargiz}.) To dig one's heels in? (\textbf{\pa{ਅੜਨਾ}}, \emph{aṛnā}.) A promise? (\textbf{\pa{ਵਾਅਦਾ}}, \emph{vāadā}.) To persuade? (\textbf{\pa{ਮਨਾਉਣਾ}}, \emph{manāuṇā}.) Absolutely? (\textbf{\pa{ਬਿਲਕੁਲ}}, \emph{bilkul}.) To admit? (\textbf{\pa{ਕਬੂਲਣਾ}}, \emph{kabūlṇā}.) A compromise? (\textbf{\pa{ਸਮਝੌਤਾ}}, \emph{samjhautā}.) Doubt? (\textbf{\pa{ਸ਼ੱਕ}}, \emph{shakk}.) Certainty? (\textbf{\pa{ਯਕੀਨ}}, \emph{yakīn}.) Trust? (\textbf{\pa{ਭਰੋਸਾ}}, \emph{bharosā}.) A guess? (\textbf{\pa{ਅੰਦਾਜ਼ਾ}}, \emph{andāzā}.) To seem right? (\textbf{\pa{ਜਚਣਾ}}, \emph{jachṇā}.) To explain? (\textbf{\pa{ਸਮਝਾਉਣਾ}}, \emph{samjhāuṇā}.) Meaning? (\textbf{\pa{ਮਤਲਬ}}, \emph{matlab}.) A reason? (\textbf{\pa{ਕਾਰਨ}}, \emph{kāran}.) Plain to understand? (\textbf{\pa{ਸਪਸ਼ਟ}}, \emph{sapshaṭ}.) To suggest? (\textbf{\pa{ਸੁਝਾਉਣਾ}}, \emph{sujhāuṇā}.) To get muddled? (\textbf{\pa{ਉਲਝਣਾ}}, \emph{ulajhṇā}.) To slip up? (\textbf{\pa{ਉੱਕਣਾ}}, \emph{ukkṇā}.) A hint? (\textbf{\pa{ਇਸ਼ਾਰਾ}}, \emph{ishārā}.) A matter? (\textbf{\pa{ਮਾਮਲਾ}}, \emph{māmlā}.) Proof? (\textbf{\pa{ਸਬੂਤ}}, \emph{sabūt}.)
\end{tcolorbox}
